% TOP_PROBLEM: {"id": 30006681, "title": "BQP Versus NP: Complete Language-Class Relation", "category_id": 15, "category": {"id": 15, "name": "computer_science", "display_name": "Computer Science", "description": "Computational complexity, algorithms, and theoretical CS.", "slug": "computer-science", "order_index": 15, "created_at": "2026-07-31T15:26:25.671Z"}, "status": "open"}
% FIRST_POSTED: null
% !TeX program = lualatex
% Compile twice: lualatex open_problems_500.tex
% All references and prose are embedded; no BibTeX database or external inputs.
\documentclass[11pt,a4paper]{article}
\usepackage[margin=24mm,headheight=14pt]{geometry}
\usepackage{fontspec}
\setmainfont{Latin Modern Roman}
\setsansfont{Latin Modern Sans}
\setmonofont{Latin Modern Mono}
\usepackage{amsmath,amssymb,mathtools}
\usepackage{unicode-math}
\setmathfont{Latin Modern Math}
\usepackage{microtype}
\usepackage{enumitem,xurl,needspace}
\usepackage[dvipsnames]{xcolor}
\usepackage{hyperref}
\hypersetup{unicode=true,colorlinks=true,linkcolor=MidnightBlue,urlcolor=MidnightBlue,
 pdftitle={Mathematical Research Notebook},pdfauthor={Source-annotated compilation},bookmarksnumbered=true}
\usepackage{bookmark}
\usepackage{tocloft}
\setlength{\cftsecnumwidth}{3em}
\cftsetpnumwidth{2.5em}
\cftsetrmarg{4em}
\usepackage{titlesec}
\titleformat{\section}{\Large\bfseries\raggedright}{\thesection}{0.7em}{}
\usepackage{fancyhdr}
\pagestyle{fancy}\fancyhf{}
\fancyhead[L]{\small\sffamily Open Problems Math | Research notebook}
\fancyhead[R]{\small 27 September 2026}
\fancyfoot[C]{\thepage}
\setlength{\parindent}{0pt}
\setlength{\parskip}{4pt plus 1pt}
\setlength{\emergencystretch}{3em}
\setcounter{tocdepth}{1}
\setcounter{secnumdepth}{1}
\setlist[itemize]{leftmargin=3em,itemsep=3pt,topsep=4pt,parsep=0pt}
\allowdisplaybreaks
\newcommand{\N}{\mathbb{N}}
\newcommand{\Z}{\mathbb{Z}}
\newcommand{\Q}{\mathbb{Q}}
\newcommand{\R}{\mathbb{R}}
\newcommand{\C}{\mathbb{C}}
\newcommand{\F}{\mathbb{F}}
\newcommand{\A}{\mathbb{A}}
\newcommand{\PP}{\mathbb P}
\newcommand{\E}{\mathbb{E}}
\newcommand{\eps}{\varepsilon}
\newcommand{\dd}{\,\mathrm d}
\newcommand{\sref}[2]{\hyperlink{source:#1:#2}{[S#2]}}
\newcommand{\eref}[2]{\hyperlink{extra:#1:#2}{[E#2]}}
% Turn the next switch off for a shorter reading copy. The quotations preserve
% every exactTarget field, including qualifications omitted from a short summary.
\newif\ifshowcatalogue
\showcataloguetrue
\newcommand{\cataloguescope}[1]{\ifshowcatalogue
 \paragraph{Exact scope in the supplied catalogue.}
 \begin{quote}\small #1\end{quote}\fi}
\usepackage{etoolbox}
\AtBeginEnvironment{itemize}{\small}
% Standalone entry; original section content is delimited for verification.
\begin{document}
\setcounter{section}{0}
%% BEGIN ORIGINAL SECTION CONTAINER
% ================================================================
% problemId: 30006681
% Canonical problem ID: 30006681
% exactTarget SHA-256: 1fdb3e1c6ae431d02de8d7efbb481b1266ce416a099474df6f498e341e16de4b
\clearpage
\section*{\texorpdfstring{BQP Versus NP: Complete Language-Class Relation}{BQP Versus NP: Complete Language-Class Relation}}\label{30006681}
\subsection{Definitions and mathematical statement}
A total language belongs to $\mathsf{BQP}$ if uniformly generated polynomial-size quantum circuits accept yes-inputs with probability at least $2/3$ and no-inputs with probability at most $1/3$. Use precisely the fixed gates, ancillas, and measurement convention in the recorded scope. The class $\mathsf{NP}$ instead uses polynomial-length classical witnesses and deterministic verification. Determine the pair of truth values
\[
 \bigl(\mathsf{NP}\subseteq\mathsf{BQP},\quad
       \mathsf{BQP}\subseteq\mathsf{NP}\bigr).
\]
A complete answer distinguishes equality, one of the two strict containments, or incomparability. Oracle access, advice, and arbitrary uncomputable gate constants are excluded.
\par\smallskip\textit{Formulation and concept references:} \sref{30006681}{1}.
\cataloguescope{Determine the truth values of BOTH NP subseteq BQP and BQP subseteq NP for unrelativized, advice-free classes of total languages over \{0,1\}. NP consists of languages with polynomial-length binary witnesses and a deterministic polynomial-time verifier. BQP consists of languages decided by polynomial-time uniformly generated polynomial-size quantum circuits (a deterministic Turing machine outputs Q\_n on 1\textasciicircum{}n), using Watrous's fixed Toffoli, Hadamard and phase-i gates with |0> ancillas and discarding, followed by a standard-basis output measurement. On every yes input acceptance is at least 2/3, and on every no input it is at most 1/3. No arbitrary real gate constants, advice or oracle are available. Establish equality, either proper containment, or incomparability.}
\subsection{Short English statement}
Which is stronger: efficient quantum decision-making or efficient classical verification, or are their powers incomparable?
%% BEGIN RESEARCH INSERTION
\subsection{Research attempt}
\paragraph{Current frontier.}
\textit{Research check: 27 September 2026.}
Watrous's Section III.2 explicitly uses the Toffoli, Hadamard, phase-$i$,
ancillary and erasure gates recorded in this entry. His survey frequently
uses promise problems, whereas the present target restricts to total
languages; that restriction is maintained below. \sref{30006681}{1}.
The 2025 random-permutation manuscript strengthens a separation relative
to certain random permutations, not an unrelativized separation.
\sref{30006681}{4}. The established black-box search lower bound also has an
oracle hypothesis. \eref{30006681}{1}. None decides either of the two requested
containments for the ordinary classes.

\paragraph{Attack route.}
For $\mathsf{NP}\subseteq\mathsf{BQP}$ one might exploit more than the
unstructured witness-search oracle. For the opposite containment one might
compress the cancellations in a quantum path sum into a classically
checkable certificate. These are different tasks; a breakthrough on one
would not answer the other. We investigate both at the exact fixed-gate
level, locating the missing structural hypothesis in each.

\paragraph{Attempt: how far black-box witness search goes.}
Let a verifier have $m$ witness bits and put $R=2^m$. Regard its accepting
set temporarily as an arbitrary Boolean oracle on $R$ addresses. Compare
the all-zero oracle with the oracle having its unique marked address at
$j$. Purify the algorithm, defer measurements, and let $q_{t,j}$ be the
probability of querying address $j$ immediately before query $t$ in the
all-zero computation. A telescoping hybrid, replacing queries in an order
that leaves the prefix equal to the all-zero computation, gives
\[
 \|\psi_j-\psi_0\|\leq2\sum_{t=1}^T\sqrt{q_{t,j}}.
\]
The suffix is unitary and does not increase the norm; a changed query
acts only on the address-$j$ subspace and has operator norm difference
at most two. Cauchy--Schwarz, followed by $\sum_jq_{t,j}=1$, yields
\[
 \sum_{j=1}^R\|\psi_j-\psi_0\|^2\leq4T\sum_{t,j}q_{t,j}=4T^2.
 \tag{16.1}
\]
For bounded-error distinction, the measured acceptance probabilities differ
by at least $1/3$ for every $j$. Their difference is at most trace distance,
which for pure states is at most their Euclidean distance. Hence
$T\geq\sqrt R/6$. This is the standard hybrid obstruction, derived here
with the quantities charged explicitly. \eref{30006681}{1}.

It rules out a universal polynomial-query search using only verifier
calls. It does \emph{not} rule out a polynomial-time algorithm that reads
the verifier's description. For example, the circuit testing $y=j$ has
a unique witness but exposes $j$ as its hardwired bit pattern. An oracle
algorithm cannot read that pattern; a circuit-input algorithm can. To
separate NP from BQP one must build a uniformly NP-verifiable family
where even description-aware quantum computation fails. Formula (16.1)
does not construct that family.

\paragraph{Attempt: exact interference accounting in the other direction.}
Fix a total-language BQP circuit $Q_n$ and input $x$. Move all erasures to
the end by retaining their unused wires, and allocate all $|0\rangle$
ancillas initially. This leaves a pure unitary computation on polynomially
many qubits, without changing the measured output probability. Let $h$
be its number of Hadamard gates. Index paths by their successive Hadamard
output bits $a\in\{0,1\}^h$. Deterministic propagation through the Toffoli
gates computes the final basis string $z(a)$ and the accumulated phase
$i^{e(a)}$, with $e(a)\in\Z/4\Z$, in polynomial time.
The Hadamard signs contribute phases $1$ or $-1$ and are included in $e$.
Thus the amplitude at $z$ and the acceptance probability are exactly
\[
 \alpha_z=2^{-h/2}\sum_{a:z(a)=z}i^{e(a)},\qquad
 p_x=2^{-h}\!\sum_{a,b}\!
 1_{z(a)=z(b)}1_{z(a)_{\rm out}=1}i^{e(a)-e(b)}.
 \tag{16.2}
\]
Conjugate pairs $(a,b)$ cancel imaginary parts. In particular
\[
 2^hp_x=G_x:=\sum_{a,b}
 1_{z(a)=z(b)}1_{z(a)_{\rm out}=1}
 \cos\!\left(\frac{\pi(e(a)-e(b))}{2}\right)
 \tag{16.3}
\]
is an integer, expressed as a sum of polynomial-time computable values
in $\{-1,0,1\}$ over $2^{2h}$ pairs. The gate normalization in (16.2)
is important: using $2^{-2h}$ would give the wrong probability.
The path-sum/classical-simulation viewpoint is established; this explicit
integer form follows directly from the fixed gates. \sref{30006681}{1}.

Although $G_x$ has at most $h+1$ bits because $0\leq G_x\leq2^h$, merely
writing down that integer is not a certificate of its correctness.
Determining a count can be difficult even when its answer has few bits.
A specific sufficient \emph{certificate lemma} would be: for every such
uniform total-language circuit family there is a deterministic
polynomial-time verifier accepting a polynomial-length certificate exactly
when
\[
 G_x>2^{h-1},
 \tag{16.4}
\]
on the inputs of that family. The BQP gap makes this equivalent to
membership in its language. Establishing this lemma for all families
would imply $\mathsf{BQP}\subseteq\mathsf{NP}$; it is not established here.
No assertion is made for arbitrary circuit descriptions with intermediate
acceptance probabilities.

\paragraph{Attempt: a tractable cancellation representation and its failure.}
A useful restricted test is to partition the path pairs in (16.3) into
polynomially many disjoint, explicitly described affine subspaces of
$\F_2^{2h}$ on each of which the summand is a known constant. Gaussian
elimination computes each cell's size, and summing its constant times
$2^{\dim\text{cell}}$ computes $G_x$ in polynomial time. If the partition
is constructed uniformly in polynomial time, this gives a P simulation,
not just an NP certificate. Given only a purported partition, however,
checking that the summand is constant on each entire cell is itself a
universal assertion. Sampling points does not furnish deterministic
soundness of a certificate.

Moreover there is no automatic low-degree closure for symbolic paths.
Retain bits $a_1,\ldots,a_r$ and reversibly compute successive prefix
products into new zero ancillas using Toffoli gates. The final bit is
$a_1a_2\cdots a_r$, whose unique multilinear polynomial over $\F_2$
has degree $r$, despite every gate having degree at most two. The circuit
is small and the function easy; it refutes bounded-degree \emph{closure},
not efficient quantum simulation in general. It also shows why simply
tracking a fixed list of pairwise correlations is not a closed method.

\paragraph{Stress test.}
The query proof concerns arbitrary oracles and cannot be substituted for
an ordinary lower bound. The path-sum formula handles phase-$i$ and
traced-out wires, but produces an exponentially long signed count, not a
polynomial certificate. A yes-certificate for a particular guessed value
must be sound on every no-input, not merely plausible numerically.
Finally, if both containments held, closure of BQP under complement would
imply $\mathsf{NP}=\mathsf{coNP}$: the complement of any NP language
would belong to BQP and hence NP. This additional consequence does not
follow from either containment alone and is a useful dependency warning.

\paragraph{Outcome.}
\textbf{Outcome: Exploratory approach with unresolved gap.}
Both directions have been investigated separately. We obtain an exact
black-box barrier, a fixed-gate signed-count formula, and a restricted
tractable representation. Neither description-aware search hardness nor
the required universal cancellation-certificate lemma is proved. Thus
neither requested truth value, and consequently no complete language-class
relation, is claimed to have been determined.
%% END RESEARCH INSERTION
\subsection{Sources}
\begin{itemize}\raggedright
\item[\textnormal{[S1]}]\hypertarget{source:30006681:1}{} John Watrous, Quantum Computational Complexity, arXiv:0804.3401v1; language questions corroborated by Scott Aaronson lecture notes.
\url{https://arxiv.org/pdf/0804.3401v1}.
{\small\textit{Catalogue formulation source.}}
\item[\textnormal{[S2]}]\hypertarget{source:30006681:2}{} Complexity Theory, Lecture 25: Quantum Computing (2) — Markus Krötzsch, January 27, 2025.
\url{https://iccl.inf.tu-dresden.de/w/images/4/4d/CT2024-Lecture-25-overlay.pdf}.
{\small\textit{Catalogue research/status reference; not a proof endorsement.}}
\item[\textnormal{[S3]}]\hypertarget{source:30006681:3}{} Introduction to Quantum Information Science lecture notes — Scott Aaronson.
\url{https://www.scottaaronson.com/qclec.pdf}.
{\small\textit{Catalogue research/status reference; not a proof endorsement.}}
\item[\textnormal{[S4]}]\hypertarget{source:30006681:4}{} Random Permutations in Computational Complexity — Hitchcock, Sekoni and Shafei (November 11, 2025).
\url{https://arxiv.org/abs/2511.08786}.
{\small\textit{Catalogue research/status reference; not a proof endorsement.}}
%% BEGIN RESEARCH REFERENCES
\item[\textnormal{[E1]}]\hypertarget{extra:30006681:1}{} Charles H. Bennett, Ethan Bernstein, Gilles Brassard and Umesh Vazirani, \emph{Strengths and Weaknesses of Quantum Computing}, SIAM Journal on Computing 26(5) (1997), 1510--1523; arXiv:quant-ph/9701001.
\url{https://arxiv.org/abs/quant-ph/9701001}.
{\small\textit{Added research source. Quantum black-box lower bounds; the hybrid calculation here keeps the oracle hypothesis explicit. Consulted 27 September 2026.}}
%% END RESEARCH REFERENCES
\end{itemize}

%% END ORIGINAL SECTION CONTAINER
\end{document}
